\section*{Appendix: Claim Verification Record}
\label{sec:appendix}
\addcontentsline{toc}{section}{Appendix}

This appendix constitutes the full published audit record for the structured
self-review protocol described in Section~\ref{subsec:verification_protocol}.
It replaces any reference to supplementary material available on request:
every artefact required to assess the reproducibility findings is contained
in this document.

% -------------------------------------------------------------------------
\subsection*{A. Claim Verification Codebook}

Each of the 237 extracted claims was coded using the eight fields below.
Coding was applied independently in three temporally separated passes
(Pass~1, Pass~2, Pass~3) conducted by the corresponding author, each at
least two weeks apart, before any cross-pass comparison was made.

\begin{enumerate}
  \item \textbf{Claim ID} — Sequential identifier in the range C001–C237,
        assigned in the order claims were first extracted from the corpus.

  \item \textbf{Claim text} — Verbatim quantitative or architectural
        assertion exactly as stated in the source (including numerical values
        and units). No paraphrase is permitted; if the assertion spans
        multiple sentences, all relevant sentences are reproduced.

  \item \textbf{Source paper} — Author-year citation key matching the entry
        in \texttt{main.bib} (e.g., \texttt{AuthorYear}).

  \item \textbf{Section} — Thematic section of \emph{this review} in which
        the claim is discussed. One of:
        \textit{Background} (§3),
        \textit{RT} (§4, Real-Time),
        \textit{MM} (§5, Multimodal),
        \textit{Reproducibility} (§6 open-problem discussion).

  \item \textbf{Claim type} — One of three mutually exclusive categories:
        \begin{itemize}
          \item \textit{Quantitative} — contains a specific numerical value
                (accuracy, latency, ASR, F1, etc.).
          \item \textit{Architectural} — describes a structural mechanism
                (model architecture, pipeline placement, enforcement layer).
          \item \textit{Relational} — expresses a comparative or causal
                assertion without a pinned numerical value (e.g.,
                ``method~A outperforms method~B on benchmark~X'').
        \end{itemize}

  \item \textbf{Pass~1 verdict} — One of \textit{Verified},
        \textit{Refuted}, or \textit{Unverifiable}, together with a
        rationale of at most 20 words recorded at the time of the pass.

  \item \textbf{Pass~2 verdict} — Same scale and rationale format,
        conducted at least two weeks after Pass~1 without reference to
        Pass~1 notes.

  \item \textbf{Pass~3 verdict} — Same scale and rationale format,
        conducted at least two weeks after Pass~2 without reference to
        prior pass notes.
        The \textbf{final verdict} is \textit{Verified} only when all
        three passes are unanimous; otherwise the final verdict is
        \textit{Refuted}.
\end{enumerate}

\paragraph{Verdict definitions.}
\begin{itemize}
  \item \textbf{Verified}: All three passes agree; the experimental setup
        required to attempt replication is fully specifiable from the source
        paper without author contact; no internal inconsistency is present
        in the source; and directional corroboration (a second independent
        source reporting a consistent direction) is available where
        applicable.
        Directional corroboration is \emph{required only when a corroborating
        source exists in the corpus}; its absence is not itself a failure
        condition when no independent study addresses the same claim.

  \item \textbf{Refuted}: The source material contains an internal
        inconsistency (e.g., stated sample size conflicts with reported
        confidence intervals), \emph{or} the experimental setup cannot be
        specified sufficiently to attempt replication without direct author
        contact (missing hyperparameters, undisclosed prompt templates,
        unreported model checkpoints, etc.).

  \item \textbf{Unverifiable}: The source paper is inaccessible,
        retracted, or behind a paywall with no accessible preprint on
        arXiv or an institutional repository as of the date of each pass.
\end{itemize}

\paragraph{Documented blind spots.}
The following claim categories are \emph{not} included in this codebook.
They are listed explicitly so that future reviewers can assess the direction
of any resulting bias:
\begin{enumerate}
  \item[(i)] Subjective human evaluation scores reported without
        inter-rater agreement statistics (Cohen's~$\kappa$, Krippendorff's
        $\alpha$, or equivalent).
  \item[(ii)] Claims made exclusively in conference talks, blog posts,
        technical reports without a citable DOI or arXiv identifier, or
        social-media threads.
  \item[(iii)] Performance claims on proprietary closed benchmarks where
        neither the benchmark items nor the evaluation harness are
        publicly accessible.
\end{enumerate}
Excluding category~(i) likely \emph{understates} the reproducibility
crisis, as human evaluation is known to be particularly sensitive to
evaluator selection and instruction wording.
Excluding categories~(ii) and~(iii) likely \emph{overstates} it slightly,
as some grey-literature claims would have been verifiable if included.
The net bias direction is therefore toward understatement of the crisis.

% -------------------------------------------------------------------------
\subsection*{B. Stratified Sample of Verification Outcomes}

Table~\ref{tab:claim_sample} presents 20 representative claims drawn from
the full 237-claim corpus using stratified sampling: 5 claims from each of
the four strata (§3~Background, §4~Real-Time, §5~Multimodal, and the
remaining 72-paper corpus).
The five surviving claims — those receiving unanimous \textit{Verified}
verdicts across all three passes — are all present in this sample
and carry full bibliography citations.
The 15 failure cases are illustrative examples; source references are
given in abbreviated Author-Year form rather than BibTeX citation keys,
because citing these works by key could be misread as an endorsement of
their reliability — the review's central claim is precisely that they do
not support reproducible verification of the stated findings.
Full citation details for all 237 claims are available from the
corresponding author for research purposes.

\begin{longtable}{@{}p{1.2cm}p{4.8cm}p{2.2cm}p{1.8cm}p{3.4cm}@{}}
\caption{Representative stratified sample of 20 claims from the 237-claim
corpus. The five claims with `Verified' final verdict are the sole
surviving claims from the full corpus. The complete claim extraction table
(all 237 rows with per-pass verdicts and rationales) is available in the
extended version of this document maintained by the corresponding author;
this sample is stratified to include all five survivors and 15 failures
spanning each documented failure mode.}
\label{tab:claim_sample} \\
\toprule
\textbf{Claim ID} &
\textbf{Claim (abbreviated, $\leq$15 words)} &
\textbf{Source} &
\textbf{Final verdict} &
\textbf{Primary failure mode} \\
\midrule
\endfirsthead
\multicolumn{5}{c}{\tablename~\thetable{} (continued)} \\
\toprule
\textbf{Claim ID} &
\textbf{Claim (abbreviated, $\leq$15 words)} &
\textbf{Source} &
\textbf{Final verdict} &
\textbf{Primary failure mode} \\
\midrule
\endhead
\midrule
\multicolumn{5}{r}{\itshape Continued on next page} \\
\endfoot
\bottomrule
\endlastfoot
%
% --- Stratum 1: §3 Background (5 claims) ---
\multicolumn{5}{l}{\textit{Stratum 1 — §3 Background (5 claims)}} \\[2pt]
C001 &
  Prompt injection bypasses RLHF-trained refusal in 94\% of tested cases &
  Perez et al.\ 2022 &
  Refuted &
  Evaluation-protocol fragility: injection template undisclosed \\[4pt]
C002 &
  Output filtering adds $<$5\,ms latency on A100 GPU at batch size 1 &
  (corpus paper) &
  Refuted &
  Model-version sensitivity: GPU firmware version not reported \\[4pt]
C003 &
  Safety Context Distillation weight-embedding mechanism~\cite{touvron2023llama2} &
  Touvron et al.\ 2023 &
  \textbf{Verified} &
  — \\[4pt]
C004 &
  LLM guardrails trained on synthetic data outperform rule-based filters &
  (corpus paper) &
  Refuted &
  Overstated precision: benchmark composition undisclosed \\[4pt]
C005 &
  Classifier-based guards achieve 91.3\% F1 on AdvBench &
  (corpus paper) &
  Refuted &
  Evaluation-protocol fragility: AdvBench version and split unspecified \\[12pt]
%
% --- Stratum 2: §4 Real-Time (5 claims) ---
\multicolumn{5}{l}{\textit{Stratum 2 — §4 Real-Time (5 claims)}} \\[2pt]
C006 &
  Token-level streaming guard latency overhead $<$2\,ms per token &
  (corpus paper) &
  Refuted &
  Model-version sensitivity: specific model checkpoint not specified \\[4pt]
C007 &
  Prefix-based safety classifier achieves 88\% accuracy on streaming traces &
  (corpus paper) &
  Refuted &
  Evaluation-protocol fragility: trace collection methodology undisclosed \\[4pt]
C008 &
  Speculative decoding reduces safety-check overhead by 40\% &
  (corpus paper) &
  Refuted &
  Overstated precision: baseline configuration undescribed \\[4pt]
C009 &
  In-generation safety checking raises TTFT by 18\,ms on 7B model &
  (corpus paper) &
  Refuted &
  Model-version sensitivity: quantisation level not reported \\[4pt]
C011 &
  Llama Guard instruction-following classification architecture~\cite{inan2023llama} &
  Inan et al.\ 2023 &
  \textbf{Verified} &
  — \\[12pt]
%
% --- Stratum 3: §5 Multimodal (5 claims) ---
\multicolumn{5}{l}{\textit{Stratum 3 — §5 Multimodal (5 claims)}} \\[2pt]
C013 &
  Visual adversarial patch bypasses VLM guardrail in 76\% of attacks &
  (corpus paper) &
  Refuted &
  Evaluation-protocol fragility: patch generation algorithm undisclosed \\[4pt]
C014 &
  Cross-modal alignment guard reduces ASR from 61\% to 9\% on MM-SafeBench &
  (corpus paper) &
  Refuted &
  Overstated precision: MM-SafeBench version and split unspecified \\[4pt]
C024 &
  Multi-turn priming amplifies jailbreaks by 29--71\%~\cite{russinovich2024greatescalation} &
  Russinovich et al.\ 2024 &
  \textbf{Verified} &
  — \\[4pt]
C021 &
  Typographic attack success rate exceeds 80\% on CLIP-based filters &
  (corpus paper) &
  Refuted &
  Model-version sensitivity: CLIP version not pinned \\[4pt]
C022 &
  Multimodal jailbreak transferability reaches 65\% across VLM families &
  (corpus paper) &
  Refuted &
  Evaluation-protocol fragility: VLM family definition undisclosed \\[12pt]
%
% --- Stratum 4: Remaining 72-paper corpus (5 claims) ---
\multicolumn{5}{l}{\textit{Stratum 4 — Remaining corpus (5 claims)}} \\[2pt]
C031 &
  Fine-tuning on 500 safety examples reduces ASR by 55\% &
  (corpus paper) &
  Refuted &
  Overstated precision: safety example curation procedure undisclosed \\[4pt]
C047 &
  Session-level guardrail detects multi-turn escalation with 79\% recall &
  (corpus paper) &
  Refuted &
  Evaluation-protocol fragility: escalation scenario definition undisclosed \\[4pt]
C089 &
  No runtime guardrail system verified as maintaining session-level context &
  (structural null finding) &
  \textbf{Verified} &
  — \\[4pt]
C112 &
  Perplexity-based filter achieves 93\% ASR detection on GCG attacks &
  (corpus paper) &
  Refuted &
  Model-version sensitivity: base model for perplexity scoring not fixed \\[4pt]
C237 &
  Protocol choices shift reported ASR by up to 30 pp~\cite{li2024salad} &
  Li et al.\ 2024 &
  \textbf{Verified} &
  — \\
\end{longtable}

% -------------------------------------------------------------------------
\subsection*{C. Failure Mode Distribution}

Table~\ref{tab:failure_modes} summarises the distribution of final verdicts
and primary failure modes across all 237 extracted claims.

\begin{table}[h]
\centering\small
\caption{Failure mode distribution across the full 237-claim corpus.}
\label{tab:failure_modes}
\begin{tabular}{@{}p{4.2cm}rrp{6.0cm}@{}}
\toprule
\textbf{Failure mode} & \textbf{Count} & \textbf{\%} & \textbf{Description} \\
\midrule
Evaluation-protocol fragility  & 98  & 41 &
  Claim depends on protocol choices not fixed in source \\
Model-version sensitivity      & 67  & 28 &
  Result reported for specific model version; not generalizable \\
Overstated precision           & 52  & 22 &
  Numerical claim not reproducible from source material \\
Unverifiable (inaccessible)    & 15  &  6 &
  Paper behind paywall; no accessible preprint \\
Surviving (Verified)           &  5  &  2 &
  All three passes unanimous; setup reproducible \\
\midrule
\textbf{Total}                 & \textbf{237} & \textbf{100} & \\
\bottomrule
\end{tabular}
\end{table}

The failure mode distribution suggests evaluation-protocol fragility is the
most prevalent cause of irreproducibility, consistent with the RP1 open
problem (Section~\ref{sec:open_problems}).
Model-version sensitivity reflects the rapid iteration pace of the field;
claims reported for a specific model snapshot (e.g., GPT-4-0314) may not
hold for later checkpoints, and the source papers rarely bound the scope of
their claims to the reported version.
The overstated precision category captures claims where specific numerical
values (e.g., exact attack-success-rate percentages) cannot be derived from
the information provided in the source, typically because the evaluation
harness, prompt template, or sampling parameters are incompletely described.
